% page 493 of the facsimile (image p-492 of the scan)
% block 170.2 x 250.6 mm ; left margin 31.8 mm
\pgc{10.160mm}{0.000mm}{
\l{\cel{55}485}
\l{}
\l{\sou{recidivar}. (\sou{netrans., pri}). Iterar krimino, delikto, kulpo pri}
\l{\cel{3}qui lu ja kondamnesis e punisesis. - DEFI.}
\l{}
\l{\sou{recipiento}. (\sou{tekn.}) Vazo adaptita ad ula aparati por recevar gasi,}
\l{\cel{3}liquida, e c. - DEFIS.}
\l{}
\l{\sou{reciproka}. I. Dicesas pri du termini, singla de qui agas a la al-}
\l{\cel{3}tra ye maniero qua equivalas olta quan lu recevas. - II. (\sou{gram.)}}
\l{\cel{3}\sou{Propozicioni reciproka}\cel{1}: Di qui singla esas nur la altra renver-}
\l{\cel{3}sita, singla di qui havas kom subjekto la atributo, e kom atri-}
\l{\cel{3}buto la subjekto di la altra. -\cel{1}\sou{Verbo reciproka}\cel{1}: Qua expresas}
\l{\cel{3}la ago da du subjekti, ica ad ita, ed ita ad ica. - III. Fondita}
\l{\cel{3}sur interkambio di agi, di sentimenti qui inter-korespondas.}
\l{\cel{3}- DEFIRS.}
\l{}
\l{\sou{recitar}. (\sou{trans., ad, koram}). Dicar voce, parole, memore, quale se}
\l{\cel{3}onu lektus. - DEFIS.}
\l{}
\l{\sou{recitativo}. I. (\sou{muziko)}. En ensemblo, solo kantala od instrumen-}
\l{\cel{3}tala. Kanto deklamata sen simetrio di mezuro e di ritmo, e ka-}
\l{\cel{3}dencata segun la dividuro di la versi e di la flexi di la}
\l{\cel{3}voco parolanta. - DEFIRS.}
\l{}
\l{\sou{reda}.I. Di qua la koloro esas olta di sango, di fairo, e c. - II.}
\l{\cel{3}(\sou{metaf.}) La\cel{1}\sou{redi}\cel{1}: la revolucionisti qui havas kom emblemo}
\l{\cel{3}standardo reda. - E.}
\l{}
\l{\sou{redaktar}. (\sou{trans.}) Kompozar segun la formo qua konvenas (libro,}
\l{\cel{3}artiklo di jurnalo o di revuo, letro, kontrato, e c.) - DEFIRS.}
\l{}
\l{\sou{redano}. (\sou{fortifik.}) Fortifiko-masivo dento-forma, ek du facii qui}
\l{\cel{3}intersekas, tale formacante angulo konvexa. - DEF.}
\l{}
\l{\sou{redemtar}.- (\sou{trans.}) I. Obtenar la liberigo po ransono. - II. Libe-}
\l{\cel{3}rigar de obligeso per pagar pekunio-quanto. - III. (\sou{metaf.})}
\l{\cel{3}(\sou{kristanismo}). Obtenar la salveso eterna po sua morto. (\sou{aludante}}
\l{\cel{3}Iesu. - EFIS.}
\l{}
\l{\sou{redingoto}. Vesto di qua la paneli envelopas la korpo. - DF.}
\l{}
\l{\sou{reduito}. Fortifikajo konstruktita dop altra por prolongar la}
\l{\cel{3}defenso-fortifikaji se la unesma esus kaptita da la enemiko.-}
\l{\cel{3}DEFIRS.}
\l{}
\l{\sou{reduktar}. (trans.). I. Retroduktar a lua situeso naturala, normala}
\l{\cel{3}(osto luxacita, ruptita). - II. Retroduktar a formo plu elemen-}
\l{\cel{3}ta (kom ex.: reduktar grani a farino, ligno a cindri, substanco}
\l{\cel{3}a pulvero, a pudro, e c.) - DEFIRS.}
\l{}
\l{\sou{reduto}. (\sou{milit-arto)}. Fortifikajo aparta, tote klozita (+kluza),}
\l{\cel{3}sen anguli konkava. - Anke : Fortifikajo flotacanta, garnisita}
\l{\cel{3}ye bastioni, uzata por transportar trupi de soldati trans flu-}
\l{\cel{3}vio o rivero. - DEFS.}
\l{}
\l{refo. (navig.) Parto di seglo quan onu rifaldis por diminutar la}
\l{\cel{3}surfaco opozata a vento. -}
}
